\documentclass[border=0pt]{standalone}
\usepackage[T1]{fontenc}
\usepackage[utf8]{inputenc}
\usepackage[spanish,es-noshorthands,es-nodecimaldot]{babel}
\usepackage{amsmath,amssymb}
\usepackage{newpxtext,newpxmath}
\usepackage{xcolor}
\usepackage{tikz}
\usetikzlibrary{arrows.meta,positioning,fit,backgrounds,calc}
\definecolor{navy}{HTML}{17243A}
\definecolor{blue}{HTML}{2A5CAA}
\definecolor{orange}{HTML}{D2691E}
\definecolor{cream}{HTML}{FCEFE3}
\definecolor{pale}{HTML}{EEF2F8}
\definecolor{text}{HTML}{1E232B}
\definecolor{muted}{HTML}{5F6873}
\definecolor{rule}{HTML}{C8D0DA}
\definecolor{sky}{HTML}{9DB9E3}
\definecolor{redact}{HTML}{3C4552}

\newcommand{\no}[1]{{\color{orange}#1}}
\setlength{\parindent}{0pt}
\setlength{\parskip}{0.45em}
\begin{document}
\hyphenpenalty=10000\exhyphenpenalty=10000
\begin{minipage}{26.783cm}\fontsize{27}{33.5}\selectfont\raggedright 
Para un estado $x=(j,a,w)$ \,---marca, edad y desgaste del coche---\, un hogar de
tipo $\tau$ obtiene por usarlo
\[ u_\tau(x)=u_{0,\tau j}+u_{1,\tau j}\,a+u_w\,w , \]
y si con la decisión $d$ se queda con el coche $h=(i,b,y)$, su utilidad del año es
\[ U_\tau(d)=u_\tau(h)-\mu_\tau\Big[\text{gastos}-\text{ingresos}\Big]. \]
El valor de cada decisión suma el futuro: si el coche se avería pasa a la edad
terminal $\bar a$; si sobrevive, envejece y su desgaste cambia,
\[ \begin{aligned} v(x;d)=U_\tau(d)+\beta\Big[\,&s(h)\,EV(i,\bar a)\\
   &+\big(1-s(h)\big)\!\int\! f(w'\mid h,r)\,EV(i,b{+}1,w')\,dw'\Big]. \end{aligned} \]
La ecuación de Bellman elige la mejor decisión, con choques $\epsilon$ de valor
extremo:
\[ V(x,\epsilon)=\max_{d\in C(x)}\Big\{v(x;d)+\epsilon(d)\Big\},\qquad
   EV(x)=\mathbb{E}_\epsilon\Big[V(x,\epsilon)\Big]. \]
\end{minipage}
\end{document}
